%%
%% test_listas.tex -- regression test for the coppe class.
%%
%% Guarda a lista de abreviaturas e siglas e a lista de simbolos:
%%
%%   * SEM numero de pagina e SEM pontilhado. A 4.1.1 descreve as duas como
%%     relacao alfabetica do termo e do seu SIGNIFICADO, e nao pede localizacao.
%%     O folio era heranca de as duas terem sido construidas com a maquina do
%%     indice remissivo. Se ele voltar, volta aqui.
%%   * a chave de ordenacao opcional, que e o que poe IoT, eMBB e uRLLC no lugar
%%     certo do alfabeto -- o makeindex ordena por codigo de caractere e sem a
%%     chave joga toda minuscula para depois de toda maiuscula.
%%   * uma descricao LONGA, que quebra em duas linhas. Era o caso feio: com o
%%     folio no fim, a quebra deixava o numero sozinho na linha seguinte.
%%
%% Nao basta compilar: as listas so aparecem depois do makeindex com o estilo
%% coppe.ist, entre a primeira e a segunda passada. O build-check.ps1 faz isso.
%%
%% NAO e gerado pelo coppe.dtx e NAO e distribuido: a suite prova que a
%% classe funciona, nao faz parte dela. Rode por tools\prova.ps1, ou por
%% tests\run-tests.ps1 para so esta pasta.
%%

\documentclass[dsc]{coppe}

\usepackage[brazilian]{babel}

\makelosymbols
\makeloabbreviations

\title{Título de Teste das Listas}
\foreigntitle{Lists Test Title}
\author{Nome}{de Teste}

\advisor{Primeiro}{Orientador}{D.Sc.}{UFRJ}
\examiner{Primeiro Examinador}{D.Sc.}{UFRJ}
\examiner{Segundo Examinador}{Ph.D.}{UFRJ}

\department{PESC}
\date{09}{2026}
\dataaprovacao{15 de setembro de 2026}

\keyword{Listas}
\keyword{Siglas}

\areaconcentracao{Engenharia de Sistemas e Computação}

\begin{document}
  \maketitle
  \frontmatter
  \begin{abstract}
    Resumo de teste.
  \end{abstract}
  \begin{foreignabstract}
    Test abstract.
  \end{foreignabstract}
  \printloabbreviations
  \printlosymbols
  \tableofcontents
  \mainmatter

  \chapter{Introdução}

  Siglas sem chave de ordenação, que saem na ordem que o \texttt{makeindex}
  quiser: \abbrev{UFRJ}{Universidade Federal do Rio de Janeiro}, \abbrev{IoT}{IoT
  com ordenação padrão}.

  Siglas com chave de ordenação, que é o que as põe no lugar certo do alfabeto:
  \abbrev[IOT]{IoT}{IoT ordenado como IOT},
  \abbrev[EMBB]{eMBB}{eMBB ordenado como EMBB},
  \abbrev[URLLC]{uRLLC}{uRLLC ordenado como URLLC}.

  Uma descrição comprida de propósito, para forçar a quebra de linha:
  \abbrev[CPGP]{CPGP}{Comissão de Programas de Pós-Graduação e Pesquisa de
  Engenharia do Instituto Alberto Luiz Coimbra, da Universidade Federal do Rio
  de Janeiro}.

  \section{Símbolos}

  Símbolos sem chave: \symbl{$\emptyset$}{Conjunto vazio},
  \symbl{$\mathbb{R}$}{Conjunto dos números reais}.

  Símbolo com chave de ordenação: \symbl[Alpha]{$\alpha$}{A letra alfa}.

  Um símbolo com descrição comprida:
  \symbl[Sigma]{$\sum_{i=1}^{n} x_i$}{Somatório dos elementos de um vetor de $n$
  posições, usado em toda parte deste documento de teste e escrito por extenso
  para forçar a quebra}.
\end{document}
%%
%%
%% End of file `tests/test_listas.tex'.
